\documentclass[12pt]{article}
\usepackage[utf8]{inputenc}
\usepackage{amsmath}
\usepackage{amssymb}
\usepackage{geometry}
\usepackage{graphicx}
\usepackage{hyperref}

\geometry{margin=1in}

\title{Optimal Seeding in Simplicial Contagion Models}
\author{Thomas Blalock}
\date{March 2026}

\begin{document}

\maketitle

\section{Introduction}
Contagion refers to the process by which a pathogen such as a virus, information, or an opinion spreads across a network. Traditional models of contagion rely on graphs where interactions are strictly dyadic, meaning that the contagion is modeled to traverse along the simple edges of the network. This approach is inadequate in the context of social contagion, where contagion through group interactions behaves fundamentally differently from contagion through pairwise interactions. For example, someone adopting a norm because other people independently convinced them to can not be equated to them adopting a norm because everyone in a participating group is an adopter. This is the difference between simple contagion and complex contagion.

To model this difference between pairwise interactions and group interactions, Iacopini et al. \cite{iacopini2019} formalize different interactions as simplices. Recalling the definition of a simplex, a k-simplex $\sigma$ is a set of k+1 vertices. A 0-simplex would be just a single node. A 1-simplex would be a pair of nodes with an edge connecting them. A 2-simplex or higher would be the set of k+1 nodes with a hyperedge connecting them in a group interaction. 

With this definition of group interactions as simplices, Iacopini et al. \cite{iacopini2019} proposed the simplicial contagion model (SCM) as a modeling framework for contagion in social contexts, unifying simple and complex contagion into a single coherent modeling framework. The SCM treats different k-simplex interactions as containing fundamentally different processes by describing contagion as a stochastic process where the infection rates depend on the dimensionality of the interaction. In simpler terms, the infection rate of a 1-simplex ($\beta$) and that of a 2-simplex ($\beta_{\Delta}$) are independent of each other. When modeled this way, group influence cannot be decomposed into independent pairwise edges.

Iacopini et al. \cite{iacopini2019} tested the SCM on several empirical and synthetic networks, yielding two primary conclusions: (i) higher-order interactions change the epidemic threshold transition from continuous to discontinuous and (ii) a bistable region emerges where both healthy and endemic states co-exist depending on the initial infection density. The SCM featuring a first-order phase transition explains well the “critical mass” observed to be required to trigger large-scale social change.

\section{Overview of the Reproduced Result}
To perform their experiments, they generated random simplicial complexes (RSC) on a scale ranging from low rescaled infectivity to high rescaled infectivity, calculated as a function of the pairwise infection rate $\beta$, the recovery rate $\mu$, and the average node degree $\langle k \rangle$ according to the formula: $\lambda = \frac{\beta \langle k \rangle}{\mu}$ where $\lambda$ represents the base rescaled infectivity. This rescaled infectivity score represents the potency of the contagion, meaning that a higher rescaled infectivity should correlate with a higher final infection density. They simulated the SCM over these synthetic networks at varying levels of 2-simplex rescaled infectivity rates ($\lambda_{\Delta}$). The results from these simulations display the discontinuous phase transition.

Given the use of RSCs, they characterized the discontinuous phase transition using the mean field approach, which provided a theoretical convergence node density based on the assumption of homogeneous mixing to check their simulations against. This yielded the following analytical solution for final node density: 
\begin{equation}
\rho^* = \frac{(\lambda_{\Delta} - \lambda) + \sqrt{(\lambda - \lambda_{\Delta})^2 - 4\lambda_{\Delta}(1-\lambda)}}{2\lambda_{\Delta}}
\label{eq:rho_star}
\end{equation}
where $\lambda_{\Delta} = \frac{\beta_{\Delta} \langle k_{\Delta} \rangle}{\mu}$ is the 2-simplex rescaled infectivity. This analytical solution correctly captures the phase transition since the discriminant is negative, the non-zero roots become complex, meaning that the endemic stationary state ceases to exist in the real domain, leaving only the absorbing state at zero.

Figure 1 plots these simulations for varying values of rescaled infectivity and $\lambda_{\Delta}$, showing a discontinuous phase transition for each set of simulations using a non-zero $\lambda_{\Delta}$. But, when $\lambda_{\Delta}$ is zero, meaning that the model is reduced to a simple contagion model without higher order interactions, the visualization displays a second order continuous phase transition.

Figure 2 plots simulations for varying initial densities of infected nodes over time with $\lambda$ fixed at 0.75 and $\lambda_{\Delta}$ fixed at 2.5. For this configuration, the theoretical initial node density threshold is $\rho_c = \frac{(\lambda_{\Delta} - \lambda) - \sqrt{(\lambda - \lambda_{\Delta})^2 - 4\lambda_{\Delta}(1-\lambda)}}{2\lambda_{\Delta}}$, plotted as a dashed line. These temporal evolutions of the infection density show the simulations starting above the threshold converging to the theoretical final infection density while the simulations with an initial infection density below the threshold converge to an infection density of zero. This shows the bistable region where the infection density is either zero or converges into a state where both infected and healthy nodes coexist.

\begin{figure}[ht]
    \centering
    \begin{minipage}{0.45\textwidth}
        \centering
        \includegraphics[width=\textwidth]{fig1a.png}
        \caption{Original Phase Transition \cite{iacopini2019}}
    \end{minipage}
    \hfill
    \begin{minipage}{0.45\textwidth}
        \centering
        \includegraphics[width=\textwidth]{fig1b.png}
        \caption{Original Bistable Temporal Evolution \cite{iacopini2019}}
    \end{minipage}
\end{figure}

\section{Criteria for Success}
We aim to replicate their conclusion regarding the SCM yielding a discontinuous phase transition and the emergence of the bistable region dependent on initial node density. In order to perform this replication, we aim to recreate their simulations over synthetic 2-simplex networks, recreating figures 1 and 2.

To recreate figure 1, we must simulate the SCM for varying rescaled infectivities and values of $\lambda_{\Delta}$. Our simulation results must feature a discontinuous phase transition that approximately coincides with the mean field approach analytical solution as described in Equation \ref{eq:rho_star}. Including the mean field approach solution validates the implementation of our random graph generation and simulation mechanics. Ensuring that our simulation results approximately match, not only the results in Iacopini et al. \cite{iacopini2019}, but also the theoretical results as calculated by the mean field approach verifies that our models are not globally biased in either direction.

To recreate figure 2, we must simulate the SCM over varying values of initial node density with a fixed rescaled infectivity, track the infection density over time, and observe the expected hysteresis loop and bistability. The results of these simulations must produce a bistable region, with simulations starting above the theoretical threshold as calculated by the mean field approach converging to the predicted final infection density and simulations below that threshold converging to zero. In these simulations, we expect the threshold to be sensitive to stochastic noise in the SCM relative to the size of the random network. This is because, on small networks, random chance may push the infection density of simulations starting near the threshold above the threshold. On larger networks, the noise will be less impactful to the overall network density, and thus the results will be less sensitive to noise.


\section{Discussion of Results}
The replication successfully reproduced the discontinuous phase transition and bistability inherent to the SCM. Intuitively, this discontinuity arises because 2-simplices require multiple active neighbors to transmit the contagion, creating a synergistic effect that only sustains itself once a critical mass (the threshold density) is reached. Furthermore, because these simulations utilized RSCs—which possess a binomial degree distribution and lack dense community clustering—the structural distinction between random and deliberate initial seeding is minimal, allowing the homogeneous mixing assumption of the mean-field approach to hold accurately. 

Our simulations, using the same simulation and graph construction hyperparameters as the original paper, are shown in figures 3 and 4. Figure 3 shows the discontinuous phase transition for each SCM simulation, and each simulation approximately matches the mean field approach analytical solution. Figure 4 shows the bistable region dependent on initial node density, with simulations starting above the density threshold converging to an endemic state and simulations starting below the density threshold converging to a healthy state. 

\begin{figure}[ht]
    \centering
    \begin{minipage}{0.45\textwidth}
        \centering
        \includegraphics[width=\textwidth]{fig2a.png}
        \caption{Replicated Phase Transition ($N=2000$)}
    \end{minipage}
    \hfill
    \begin{minipage}{0.45\textwidth}
        \centering
        \includegraphics[width=\textwidth]{fig2b.png}
        \caption{Replicated Bistable Region ($N=2000$)}
    \end{minipage}
\end{figure}

Notably, two simulations with an initial node densities below the threshold, converged to the endemic state. One explanation is that this exception was due to random noise since the network was small (2,000 nodes). Another explanation is that the theoretical threshold is too high since its derivation assumed homogeneous mixing. To test these explanation, we reran the simulation over larger random networks with 8,000 nodes. If the anomaly was due to random chance, then the scale would reduce the impact of noise and the results would adhere to the theoretical threshold. The results of these simulations, shown in figure 5, did not coincide with the theoretical results. Thus, we concluded that the real threshold for our random graphs was lower than the theoretical threshold derived via the mean field approach.

\begin{figure}[ht]
    \centering
    \includegraphics[width=0.7\textwidth]{fig3.png}
    \caption{Noise Reduction and Theoretical Convergence ($N=8000$)}
\end{figure}

\section{Reflection}
The process of replicating the results of Iacopini et al. \cite{iacopini2019} taught me to have redundant tests that my experiments are working correctly. I had an initial false start regarding the construction of the random graphs where I had neglected to add the simple edges that corresponded to the 2-simplex edges. My initial batch of simulations still featured the discontinuous phase transition. So, if that were my only validating check, then I might have missed the oversight in my code. However, visualising the results against the mean field approach analytical solution made the nature of the problem obvious. The final infection densities of my simulations were globally negatively biased. The redundant checks helped me catch a vital error. The observation also narrowed the location error substantially.

The act of building and running the experiments also made the limitations of their conclusions obvious, specifically regarding the use of RSCs and random seeding. For me, a lot of the time, it is difficult to determine where there is opportunity for incremental improvements on existing methods. Instead, if a methodology does not exactly meet my needs, I tend to attempt to create my own methodologies from scratch. Recreating these experiments made it obvious to me how I could incrementally build upon their framework to solve problems relevant to my projects. I chose to replicate this paper because I figured it would be useful when building an algorithm for my app deciding to which users I would send notifications regarding a new event. Before replicating the results of this paper, I envisioned jumping straight to optimizing the placement of seed nodes, which would have caused me to skip steps in the research process. This replication made it more obvious to me that I should first seek to understand the properties of the SCM over networks that better resemble social networks than RSCs, in my context with community structure and exponentially distributed degree. My intuition is that random seeding within communities or clusters biased by degree would dramatically lower the initial infection density threshold, providing a much better baseline for my eventual optimization algorithm. Additionally, it would be useful to understand how these models behave when noise or personalized beta infection rates are introduced to the system. Would noise affect the endemic threshold consistently in either direction, randomly, or not at all? Forcing myself to actually replicate the paper rather than just read it gave me a more concrete direction for future research.

\section*{Code Availability}
The source code for the simulations, graph generation, and data analysis used in this replication is available on GitHub at: \url{https://github.com/ThomasBlalock/Optimal-Seeding-in-Simplicial-Contagion-Models}.

\begin{thebibliography}{9}

\bibitem{iacopini2019} 
Iacopini, I., Petri, G., Barrat, A., \& Latora, V. (2019). 
Simplicial models of social contagion. 
\textit{Nature communications}, 10(1), 2485.

\end{thebibliography}

\end{document}